\renewcommand{\thesection}{S\arabic{section}}
\renewcommand{\thefigure}{S\arabic{figure}}
\renewcommand{\thetable}{S\arabic{table}}
\setcounter{section}{0}\setcounter{figure}{0}\setcounter{table}{0}

\section{Solver verification and the gradient oracle}
\label{si:solver}

The projection defined in the method is solved by a one-dimensional monotone root
find on the multiplier. Two things are verified rather than assumed. First, the solve
%% results/CP2_hessian_speedup_sweep.json, results/CP6_multiconstraint.json
converges to the constraint to machine precision wherever the target is attainable, and
reports non-convergence rather than returning a near-miss where it is not; the
non-convergence flag is what the fallback path keys on. Second, the analytic Jacobian
agrees with finite differences and with autograd through an unrolled solver.

For the multi-constraint form, convergence alone is insufficient: a converged solve on a
near-singular system is not trustworthy. We therefore report the conditioning of the
%% results/CP6_multiconstraint.json (conditioning report, vocabulary sweep)
constraint matrix alongside every joint solve, and sweep vocabulary size rather than
testing a single configuration, because a single grid at one vocabulary would have read as
``constraints always compose''.

\section{Overhead and timing}
\label{si:overhead}

%% results/A2_training_summary.json (wall_clock_ratio_A3_over_A2, hessian_speedup_vs_batch)
Training through the projection costs $6.0\times$ the wall clock of training with the
mask alone, on the synthetic task at matched steps. The backward pass is a rank-one
correction rather than a linear solve, and batching the Hessian-vector products in
single precision gives a speedup ranging from $1.33\times$ at the smallest batch to
$2.64\times$ at the largest. At inference the layer adds one root find per step.

\section{Grammar mask and parse validity}
\label{si:parse}

%% results/WS2_parse_validity.json (arms.masked, mask_freedom, claim_scope)
The character-level grammar makes every sampled string parse: the syntactic parse rate
is $1.00$ under an independent parser. It is not a validity oracle: only $47\%$ of
those strings build into a structure, the remainder failing on semantic grounds such as
overlapping sites or unphysical stoichiometries, which parse without complaint. The mask
leaves real freedom: a mean of $4.14$ allowed characters per step, though $70.6\%$ of
steps are fully forced by the grammar.

\section{How the overhead varies with constraint strength}
\label{si:frontier}

%% results/WS1_steering_frontier.json (summary.*)
The surrogate's excess divergence over the sequence-level optimum is close to fixed in
absolute terms while the optimum's own cost varies over a $129$-fold range as the
constraint pulls harder. The relative waste therefore falls from $97.8\%$ at the weakest
steering to $32.4\%$ at the strongest: the surrogate is least efficient exactly where the
constraint barely matters, and most efficient where it binds. \Cref{si:solver} covers the
solve itself and \cref{si:overhead} its wall-clock cost.

\section{Energy screen calibration}
\label{si:energy}

%% results/WS3_energy_calibration.json (surrogate, held_out_agreement)
Formation energies are estimated with CHGNet~\citep{deng2023chgnet} at
version v0.3.0~\citep{chgnet_v030} rather than density functional theory. On $1{,}456$ held-out structures the surrogate agrees with its reference at
$0.0210$~\unit{\electronvolt\per atom} mean absolute error and Pearson $r = 0.9996$.
This is a surrogate screen and is labelled as such wherever it is used; no claim in the
main text depends on a DFT-quality energy.

\section{Statistical floor of the difference test}
\label{si:floor}

%% results/WS6_best_of_n.json and the enumeration in scripts/crystal/equivalence_test.py
The claim that a paired Wilcoxon test at $n = 7$ cannot return $p < 0.0156$ is verified
by exhaustive enumeration over all $2^{7}$ sign patterns rather than by appeal to a
table. This is why the equivalence procedure, rather than the
difference test, carries the retargeting claim.

\section{The full equivalence appendix}
\label{si:equivalence}

%% preregistration/WS4_equivalence_margin.md (the binding margin and decision rule)
%% results/WS4_equivalence.json (tost_p, bootstrap interval, margin_matches_stated_rule)
The margin is defined before the test, in a pre-registration file that the analysis
script reads; the script fails if the file is absent or if its value disagrees with the
stated rule. The rule is that the margin equals the retargeted arm's mean within-target
seed standard deviation. Two one-sided tests give $p = 0.0023$ at that margin of
$0.189$, and the secondary pre-registered condition, still $p < 0.05$ at a $0.10$
margin, also holds. A bootstrap interval is reported alongside the parametric one and
the wider of the two is what the main text quotes.

%% results/WS4_retargeting_seed_replication.json (summary.*)
The difference test is reported for completeness and interpreted as uninformative: the
paired Wilcoxon over $7$ targets returns $p = 1.0$, but the smallest attainable
two-sided $p$ at that sample size is $2/2^{7} = 0.0156$ under any data whatsoever, so a
large value carries almost no information.

%% results/WS4_retargeting_seed_replication.json (per_seed, summary)
The table below gives the three replication seeds behind both readings, and it is
included because the aggregate hides the reason to distrust one of them. Each row is a
separate training run of the projection arm at the same target, read first without the layer
to give its backbone and then with the layer across the whole sub-ceiling family. The
backbone column is what makes seed~0 unusable as evidence: at
$22.769$ it sits far above the other two, which land at
$8.381$ and $8.692$, so its lower
sub-ceiling error of $0.4406$ reflects a backbone that happened
to propose high-density structures rather than a layer that steered better. Averaging the
three gives $0.6497$ against the frontier's
$0.6620$, which is the comparison the main text reports, and
the spread of $0.1506$ across seeds is what the pre-registered
margin is derived from. All three seeds exceed the published ceiling of
$10.405$, so the excursion is a matter of degree rather than of kind.

\begin{table}[t]
  \centering
  %% results/WS4_retargeting_seed_replication.json and results/WS4_zero_shot_retargeting.json
  \caption{\textbf{Per-seed retargeting summary.} Densities in
    \unit{\gram\per\centi\metre\cubed}. Seed 0's backbone is an excursion, which is why
    its lower error is not the reading to take; the published ceiling is $10.405$.
    Structural validity is $1.0$ in every cell of every seed.}
  \label{si:tab:perseed}
  \small
  \begin{tabular*}{\textwidth}{@{\extracolsep{\fill}}@{}l S[table-format=2.3] S[table-format=2.3] S[table-format=1.4]@{}}
    \toprule
    \kilth{Seed} & {\kilth{Backbone mean}} & {\kilth{Max achieved}} &
    {\kilth{Sub-ceiling MAE}} \\
    \midrule
    0 & 22.769 & 15.495 & 0.4406 \\
    1 &  8.381 & 10.439 & 0.7889 \\
    2 &  8.692 & 10.596 & 0.7197 \\
    \midrule
    Mean of seeds        & {--} & {--} & 0.6497 \\
    Frontier, 7 runs     & {--} & {--} & 0.6620 \\
    \bottomrule
  \end{tabular*}
\end{table}

\section{Two-constraint grid, per cell}
\label{si:grid}

%% results/WS5_two_constraint_grid.json (summary.fixed_valence_columns)
\Cref{si:tab:grid} gives every cell of the two-constraint grid at a shared valence
target, each against its own single-constraint control. The control is necessary: the
achieved valence density rises with the density target even when no valence constraint
is applied, so an absolute level would show apparent steering where there is none. Only
the deviation from control is the second constraint's effect.

\begin{table}[t]
  \centering
  %% results/WS5_two_constraint_grid.json (summary.fixed_valence_columns[*].rows)
  \caption{\textbf{Direction of the second constraint's effect.} Cells at a shared
    valence target across density rows. ``Asked'' is the direction the target sits
    relative to that cell's control. Lowering the quantity succeeds in both cells that
    request it; raising it succeeds in two of five.}
  \label{si:tab:grid}
  \small
  \begin{tabular*}{\textwidth}{@{\extracolsep{\fill}}@{}S[table-format=1.3] S[table-format=2.1] S[table-format=1.4]
                    S[table-format=1.4] S[table-format=+1.4] l c@{}}
    \toprule
    {\kilth{Valence}} & {\kilth{Density}} & {\kilth{Achieved}} & {\kilth{Control}} &
    {\kilth{Deviation}} & \kilth{Asked} & \kilth{Correct} \\
    \midrule
    0.282 &  3.0 & 0.2447 & 0.2424 & +0.0023 & up   & yes \\
    0.282 &  5.0 & 0.3062 & 0.3141 & -0.0079 & down & yes \\
    0.282 &  8.0 & 0.3251 & 0.3956 & -0.0705 & down & yes \\
    0.484 &  3.0 & 0.2603 & 0.2424 & +0.0179 & up   & yes \\
    0.484 &  8.0 & 0.3482 & 0.3956 & -0.0474 & up   & no  \\
    0.686 &  3.0 & 0.2353 & 0.2424 & -0.0071 & up   & no  \\
    0.686 &  8.0 & 0.3707 & 0.3956 & -0.0248 & up   & no  \\
    \bottomrule
  \end{tabular*}
\end{table}

\section{The boundary of the expressible class}
\label{si:expressible}

%% results/CP7_second_channel.json (charge_is_not_linear, channel_statistics)
%% results/CP9_serialization_expressibility.json (rows, verdict)
Charge neutrality cannot be written as a per-step linear expectation over this
vocabulary. The screen that defines it quantifies existentially over oxidation-state
assignments; most elements in the split are multivalent; and collapsing to a single
charge per element destroys the predicate rather than approximating it. This is a
property of the constraint, not a gap in the implementation, and it bounds the class the
layer addresses. Valence electron density was adopted in its place because it is
computable for every element present and only partly correlated with mass density
($r = 0.725$), so a joint solve is not one constraint written twice.

A second boundary is the serialization. Over $393$ structures and $1{,}217$ element
positions, a cell-first serialization makes the density budget determinable at $100\%$
of positions and a composition-first serialization at $0\%$.

\section{Defects found, and how}
\label{si:defects}

Three classes of error were caught during this work by mechanisms worth reporting, since
each would have produced a plausible-looking wrong number.

A budget generalization initially read cell-parameter token payloads as physical values
when they are bin indices, which inflated volumes while still producing budgets that
looked reasonable; it was caught by an equivalence check against the original
single-constraint path rather than by inspection. A charge-variance comparison was made
against the unconstrained arm when the claim required the sequence-level optimum:
a true statement answering a different question, caught in review, and the corrected
figure is stronger than the original. A crossing count in an earlier draft was a literal
in the document-generating script rather than a value computed from the data, and
survived a verification pass that checked every number in the same table; it is now
recomputed and asserted by a test.

The general lesson we would pass on is that verifying numbers is not the same as
verifying prose. A test that recomputes a table's values will not catch the sentence
describing them.
